%-------------------------------------------------------------------------------
%    ABBREVIATIONS
%-------------------------------------------------------------------------------

\begin{abbreviations}{ll} % Товчлолын жагсаалт оруулах (хоёр баганатай хүснэгт)
\addchaptertocentry{\abbrevname}

\textbf{AI} & \textbf{A}rtificial \textbf{I}ntelligence (Хиймэл оюун ухаан)\\
\textbf{DAW} & \textbf{D}igital \textbf{A}udio \textbf{W}orkstation (Тоон аудио ажлын станц)\\
\textbf{CNN} & \textbf{C}onvolutinal \textbf{N}eural \textbf{N}etworks (Хувиргалтын мэдрэлийн сүлжээ)\\
\textbf{MusicVAE} & \textbf{Music} \textbf{V}ariational \textbf{A}uto\textbf{e}ncoder\\
\textbf{LLM} & \textbf{L}arge \textbf{L}anguage \textbf{M}odel (Том хэлний загвар)\\
\textbf{RAG} & \textbf{R}etrieval-\textbf{A}ugmented \textbf{G}eneration (Мэдээлэл хайлтаар баяжуулсан үүсгэлт)\\
\textbf{EQ} & \textbf{E}qualizer (Экволайзер)\\
\textbf{MIDI} & \textbf{M}usical \textbf{I}nstrument \textbf{D}igital \textbf{I}nterface\\
\textbf{DSP} & \textbf{D}igital \textbf{S}ignal \textbf{P}rocessing (Тоон дохио боловсруулалт)\\
\textbf{VST} & \textbf{V}irtual \textbf{S}tudio \textbf{T}echnology (Виртуал студи технологи)\\
\textbf{UML} & \textbf{U}nified \textbf{M}odelling \textbf{L}anguage (Загварчлалын нэгдсэн хэл)\\
\textbf{GUI} & \textbf{G}raphical \textbf{U}ser \textbf{I}nterface (Хэрэглэгчийн график интерфейс)\\
\textbf{FLP} & \textbf{FL} Studio \textbf{P}roject (FL Studio төслийн файл)\\

\end{abbreviations}